\scanpage{104}
\mspage{- 115 -}
toute sous-algèbre de type fini de $C\inserted{[}T\inserted{]}$ vérifie la condition des chaînes, il en est de même de toute sous-algèbre de type fini de $C_{\uline{\unsure{r|p}}}\inserted{[}T\inserted{]}$, ce qui achève la démonstration.

\margin{$\boxed{\begin{gathered}\dim\uline O_x\leq\dim\uline O_y+\dim(\unsure{\uline O_x\otimes k(y)})-\delta(x)\\
\delta(x)=\deleted{\unsure{\deg}}\inserted{\dim_{\unsure{x}}}f^{-1}f(x)-e\end{gathered}}$}

\uline{Corollaire 1 .- Soient $Y$ un préschéma irréductible localement noethérien}, $X$ \uline{un préschéma irréductible}, $f:X\inserted{\to}Y$ \uline{un morphisme} \deleted{\illegible{1}} \uline{dominant de type fini. Soit} $\inserted{\xi}$ (resp.\ $\inserted{\eta}$) \uline{le point générique de} $X$ (resp. $Y$) \deleted{\unsure{et}}, \uline{et soit} $e=\dim(f^{-1}(\inserted{\eta}))=\deg.\operatorname{tr}_{\inserted{\uline{\kappa}(\eta)}}\inserted{\uline{\kappa}(\xi)}$ (``dimension de la fibre générique)''). \uline{Pour tout} $x\inserted{\in}X$, \uline{on a, en posant} $y=f(x)$
\[
(13)\qquad e+\operatorname{codim}_Y\overline{\inserted{\{}y\inserted{\}}}\inserted{\geq}\deg.\operatorname{tr}_{\inserted{\uline{\kappa}(y)}}\inserted{\uline{\kappa}(x)}+\operatorname{codim}_X\overline{\inserted{\{}x\inserted{\}}}.
\]
\margin{formule hybride, donner plutôt \uline{deux} formules l'une en termes d'alg. comm., l'autre en termes ``géom.''}

\uline{En outre, les deux membres sont égaux lorsque} $\uline{O}_y$ \uline{est un quotient d'anneau local régulier}.
\margin{\deleted{\unsure{si}} cond. des chaînes \unsure{stricte}}
\margin{\uuline{Variante}\quad\uuline{en vue}\quad$\left\{\begin{gathered}\dim\uline O_x\leq\dim\uline O_y+\dim\uline O_x\otimes k(\unsure{y|x})\\
-[\dim_xf^{-1}f(x)-\unsure{e}]\end{gathered}\right.$}

Comme la question est locale, on peut supposer $X$ et $Y$ affines\inserted{(et réduits)}. \deleted{\unsure{En effet} \illegible{15} \unsure{pour} \illegible{8}}\inserted{soient} $X=\operatorname{Spec}(B)$, $Y=\operatorname{Spec}(A)$, où $B$ est une algèbre intègre de type fini sur $A$; comme $f$ est dominant, $A$ s'identifie à une sous-algèbre de $B$ (§ 1,n°2,cor.4 de la prop.5), donc $\uline{O}_y$ à une sous-algèbre de $\uline{O}_x$. Comme $\inserted{\uline{\kappa}(\eta)}$ et $\inserted{\uline{\kappa}(\xi)}$ sont les corps de fractions de $A$ et $B$ respectivement, donc de \deleted{$O_x$} $\uline{O}_y$ et $\uline{O}_x$ respectivement, la formule (13) résulte de (12) appliquée aux anneaux $\uline{O}_y$ et $\uline{O}_x$.
\margin{\deleted{Si $\unsure{Z}\subset X$ fermé irréductible, alors $\dim_{\unsure{x}}f^{-1}f(x)$} \unsure{n'énoncer} \illegible{3}}
\margin{\figure{\begin{tikzpicture}\draw[->] (2,2)--(0,0);\end{tikzpicture}}}

On notera d'ailleurs que $\deg.\operatorname{tr}_{\inserted{\uline{\kappa}(y)}}\inserted{\uline{\kappa}(x)}$ n'est autre que la dimension de $f^{-1}(y)\inserted{\cap}\overline{\inserted{\{}x\inserted{\}}}$, comme il résulte de la prop.7 du n°2, appliquée à $f^{-1}(y)$ considéré comme préschéma algébrique sur $\inserted{\uline{\kappa}(y)}$. En particulier:
\margin{\deleted{\unsure{dim}}}

\uline{Corollaire 2} .- \deleted{\unsure{Si}} \uline{Les hypothèses du cor.1 étant vérifiées, supposons en outre que $x$ soit fermé dans $f^{-1}(y)$ et que $\uline{O}_y$ soit un quotient d'anneau régulier. Alors}
\margin{\deleted{\unsure{si}} cond. des chaînes \unsure{stricte}}
\[
(14)\qquad\operatorname{codim}_X\overline{\inserted{\{}x\inserted{\}}}=e+\operatorname{codim}_Y\overline{\inserted{\{}y\inserted{\}}}.
\]
\uline{Corollaire 3 .- Les hypothèses du cor.1 étant vérifiées, on a}
\margin{103}
